\paragraph{背景与动机.}

激光超声（LUS）作为一种新兴的无损检测（NDI/NDT）技术，相比传统压电式超声方法具有独特优势。
与接触式传感器不同，LUS通过激光诱导的热弹性效应实现超声波的非接触式产生和探测~\cite{laser_us_ndt}。这种非接触特性无需耦合介质，使其能够对粗糙、高温或难以接近的表面进行检测，而这些条件下传统方法往往不可行。
此外，LUS系统产生宽带声发射，具有宽频带特性，有利于实现多种材料的高分辨率缺陷检测~\cite{laser_us_broadband}。快速光栅扫描或单点测量的能力也使LUS特别适用于吞吐量至关重要的工业检测流水线~\cite{laser_us_scanning}。这些优势推动了LUS在航空航天、制造和基础设施监测领域的广泛应用。

尽管具有理论潜力，LUS的实际应用面临一个根本性挑战：采集的信号信噪比（SNR）严重下降。
激光产生的超声振幅受激光能量约束和材料吸收特性的固有限制，而环境电磁噪声、激光源波动和环境振动进一步污染了测量波形~\cite{low_snr_challenge}。这种噪声退化在现场条件下或检测低声学阻抗材料时尤为显著，信号振幅可能比噪声底限低数个数量级。因此，缺乏有效的去噪预处理，可靠的信号解释和缺陷表征将变得困难。

传统数字信号处理（DSP）技术已广泛应用于LUS去噪，包括维纳滤波、小波阈值化和块匹配3D（BM3D）滤波~\cite{wiener_filter,wavelet_denoise,bm3d_denoise}。虽然这些方法在中等噪声条件下提供了改善，但当SNR降至实际阈值以下时往往不足。维纳滤波器依赖于难以在现实场景中估计的精确噪声统计量。
小波方法虽然在稀疏表示方面有效，但在抑制非高斯噪声伪影时难以保留精细信号特征。BM3D尽管在自然图像上表现出色，但未能充分利用超声波形特有的时间结构。
这些局限性促使了数据驱动的学习方法的开发，这些方法能够自适应地捕获LUS信号中固有的复杂噪声模式和信号相关性。